\chapter{Introduzione}

\section{Non-Playable Character:  contesto e ruolo}

I personaggi non giocanti (NPC) sono elementi fondamentali nel mondo videoludico. Sono parte integrante di un mondo interattivo e contribuiscono alla credibilità del mondo stesso. Ogni videogioco costruisce l'architettura dei propri NPC in base all'obiettivo del gioco: se parliamo di un gioco prettamente d'azione dove gli NPC rappresentano minacce per il giocatore, come per esempio nella maggior parte dei giochi della serie Super Mario, non è necessario che essi si comportino in maniera "umana", replicando tratti e azioni umane, ma devono semplicemente seguire il giocatore per attaccarlo. In altri videogiochi invece, come The Sims o Tomodachi Life, dove si simula la quotidianità della vita di un essere umano, gli NPC sono il fulcro del design dei giochi: essi devono essere in grado di agire tra di loro e nell'ambiente circostante, reagire a stimoli fisiologici e così via. In pratica, si cerca di simulare il comportamento umano.
La mente umana, d'altro canto, è molto complessa, e rappresentare i comportamenti che il pensiero genera è altrettanto complicato. Le architetture proposte tradizionalmente sono statiche: si basano su macchine a stati finiti in cui il personaggio compie azioni predefinite in risposta a condizioni predefinite.

\section{Motivazioni e Obiettivi}



